% Multimodal Gaussian Splatting SLAM: A Survey of LiDAR, ToF, and Monocular-Depth-Enhanced Reconstruction
% Compile: pdflatex main && bibtex main && pdflatex main && pdflatex main
\documentclass[10pt,a4paper]{article}

\usepackage[utf8]{inputenc}
\usepackage[T1]{fontenc}
\usepackage[margin=1in]{geometry}
\usepackage{microtype}
\usepackage{amsmath,amssymb,amsfonts}
\usepackage{mathtools}
\usepackage{booktabs}
\usepackage{graphicx}
\usepackage{xcolor}
\usepackage{hyperref}
\usepackage{cleveref}
\usepackage{enumitem}
\setlist[itemize]{leftmargin=*,nosep}
\setlist[enumerate]{leftmargin=*,nosep}
\usepackage{caption}
\usepackage{subcaption}
\usepackage{multirow}
\usepackage{makecell}

\hypersetup{colorlinks=true,linkcolor=blue!60!black,citecolor=green!50!black,urlcolor=blue!60!black}

% Custom commands
\newcommand{\SEthree}{SE(3)}
\newcommand{\Simthree}{\mathrm{Sim}(3)}
\newcommand{\Twb}{T_{W\leftarrow B}}

\title{Multimodal Gaussian Splatting SLAM: \\
       A Survey of LiDAR, Time-of-Flight, and \\
       Monocular-Depth-Enhanced Real-Time Reconstruction}

\author{Arthur Li Rui \\
  \texttt{arthurlirui@example.com} \\
  Gaussian-LIC Project
}

\date{October 12, 2026}

\begin{document}
\maketitle

\begin{abstract}
3D Gaussian Splatting (3DGS) has rapidly emerged as the dominant scene
representation for photo-realistic SLAM, supplanting NeRF-based volumetric
approaches with its explicit, differentiable, and rasterization-efficient
primitives.  However, the quality of Gaussian-based SLAM is fundamentally
constrained by the \emph{depth source} used to initialize and supervise the
Gaussian map.  This survey systematically reviews the state of the art along
three complementary axes of multimodal depth provision: (1)~\textbf{LiDAR +
RGB + IMU} fusion, where precise metric geometry from LiDAR anchors the
Gaussian map while the camera and IMU provide appearance and motion priors;
(2)~\textbf{iToF Depth + RGB + IMU} fusion, where sparse indirect
Time-of-Flight depth is densified and fused with photometric cues for
lightweight, low-power mapping; and (3)~\textbf{Monocular Depth + RGB + IMU}
fusion, where learned depth priors or visual-inertial structure-from-motion
replace active depth sensing entirely.  For each axis, we analyze the pipeline
architecture (front-end odometry, depth completion, Gaussian initialization,
and online optimization), discuss open-source code availability, and identify
the key trade-offs among accuracy, robustness, real-time performance, and
hardware cost.  We further provide a cross-category comparison and outline
open challenges and future directions for multimodal Gaussian Splatting SLAM.
\end{abstract}

%======================================================================
\section{Introduction}
\label{sec:intro}
Simultaneous Localization and Mapping (SLAM) is a foundational capability
for autonomous robots, augmented reality (AR), and autonomous navigation.
Traditional SLAM systems---whether filter-based~\cite{qin2018vinsmono,
xu2022fastlio2} or graph-based~\cite{campos2021orbslam3}---produce sparse
or semi-dense maps that suffice for localization but lack the photo-realistic
detail required for immersive visualization, semantic understanding, or
high-fidelity simulation.  The introduction of Neural Radiance Fields
(NeRF)~\cite{mildenhall2020nerf} and, more critically, 3D Gaussian
Splatting~\cite{kerbl2023_3dgs} has catalyzed a paradigm shift: SLAM systems
now aim to simultaneously estimate the camera trajectory and reconstruct a
\emph{photo-realistic, dense, and renderable} 3D map of the environment.

3DGS represents scenes as sets of anisotropic 3D Gaussian primitives with
position, covariance, opacity, and spherical-harmonic color coefficients.
Unlike NeRF's costly volumetric ray-marching, 3DGS employs fast
tile-based rasterization, achieving real-time rendering at over 100~FPS on
consumer GPUs.  This efficiency, combined with differentiable optimization,
makes 3DGS the ideal map representation for online SLAM.

However, the quality of a Gaussian-based SLAM map is fundamentally limited by
the \emph{depth source} that seeds and supervises the Gaussians.  Early 3DGS
SLAM systems~\cite{yan2024gsslam,keetha2024_splatam,matsuki2024monogs}
operated with RGB-D cameras indoors, where dense active depth is readily
available.  When deployed in large-scale outdoor environments, these systems
degrade under violent ego-motion, varying illumination, textureless regions,
and---most critically---the limited range of consumer RGB-D sensors.

Three complementary strategies have emerged to provide the geometric priors
needed for robust Gaussian SLAM beyond the indoor RGB-D setting:

\begin{enumerate}
    \item \textbf{LiDAR + RGB + IMU fusion} (Section~\ref{sec:lidar}):
    Mechanically or solid-state LiDAR provides precise, long-range, but
    sparse metric geometry.  The camera contributes rich texture, and the
    IMU bridges motion gaps.  This configuration, well-established in
    classical LIO/LIVO
    systems~\cite{xu2022fastlio2,zheng2024_fastlivo2,lin2024_r3live},
    has been extended with Gaussian Splatting mapping by a rapidly growing
    family of works~\cite{lang2026_gaussianlic2,lang2025_gaussianlic,
    hong2025_gslivo,xie2024_gslivm,zhao2025_lvigs,li2026_rmgs,
    tak2026_realtimelidar,he2026lvgsslam,shi2026_litgs}.

    \item \textbf{iToF Depth + RGB + IMU fusion} (Section~\ref{sec:itof}):
    Indirect Time-of-Flight (iToF) sensors offer low-power, compact active
    depth at the cost of extreme sparsity (as few as $8\times8$ points per
    frame) and environmental noise sensitivity.  Recent
    work~\cite{conti2025_tofsplatting} demonstrates that even such sparse
    depth, when fused with multi-view geometry and monocular cues within a
    3DGS framework, can achieve competitive SLAM performance---a crucial
    step toward mobile and AR/VR deployment.

    \item \textbf{Monocular Depth + RGB + IMU fusion} (Section~\ref{sec:monodepth}):
    When no active depth sensor is available, learned monocular depth
    priors~\cite{yang2024depthanythingv2} or feed-forward 3D foundation
    models~\cite{wang2025_vggt} can substitute.  The IMU provides metric
    scale and motion continuity.  This axis spans pure visual-inertial
    Gaussian SLAM~\cite{zhu2026_vigssslam,zhang2026_hislam2,
    seo2026_gaussianflow,zhou2026_geogs,zheng2025_wildgssslam,
    han2026_vggtgssslam,zhang2026_vimgs,deambrogi2026_cgssslam,
    sun2024_mm3dgs,zhu2026_mygosplat,kruzhkov2026_monoemgs} and is the
    most hardware-accessible but algorithmically challenging configuration.
\end{enumerate}

\textbf{Scope and methodology.}  This survey covers papers published or
posted as preprints through October~2026, identified through systematic web
searches using the Doubao search API with keywords combining ``Gaussian
Splatting,'' ``SLAM,'' ``LiDAR,'' ``ToF,'' ``monocular depth,'' ``visual
inertial,'' and related terms.  For each paper, we downloaded the arXiv PDF
(when available), read the full method section, and searched for open-source
code repositories.  Papers already present in the project's related-work
bibliography (130+ entries accumulated over 30+ survey addenda) were
cross-referenced to avoid duplication.  We focus on systems that perform
\emph{online} or \emph{real-time} SLAM with 3DGS as the primary map
representation, excluding offline 3DGS reconstruction and non-SLAM novel-view
synthesis.

\textbf{Contributions.}  This survey provides:
\begin{itemize}
    \item A structured taxonomy of multimodal Gaussian Splatting SLAM along
    three depth-source axes (LiDAR, iToF, monocular depth).
    \item Per-paper analysis of architecture, depth fusion strategy, code
    availability, and key results for 25+ representative systems.
    \item A cross-category comparison identifying the accuracy--cost--robustness
    trade-off landscape.
    \item A discussion of open challenges: depth completion under degeneracy,
    scale consistency, memory-bounded mapping, and real-time deployment on
    embedded platforms.
\end{itemize}


%======================================================================
\section{Background: 3D Gaussian Splatting for SLAM}
\label{sec:background}
\subsection{3D Gaussian Splatting Primer}

3D Gaussian Splatting~\cite{kerbl2023_3dgs} represents a 3D scene as a set
of $N$ Gaussian primitives, each parameterized by:
\begin{itemize}
    \item Position $\boldsymbol{\mu}_i \in \mathbb{R}^3$ (center)
    \item Covariance $\boldsymbol{\Sigma}_i = \boldsymbol{R}_i
    \boldsymbol{S}_i \boldsymbol{S}_i^T \boldsymbol{R}_i^T$ (decomposed into
    rotation $\boldsymbol{R}_i$ and scale $\boldsymbol{S}_i$ for positive
    semi-definiteness)
    \item Opacity $\alpha_i \in [0,1]$
    \item Color via spherical harmonics (SH) coefficients up to degree~3
\end{itemize}
Rendering is performed via differentiable tile-based rasterization: each
pixel's color is an alpha-composited blend of projected Gaussians,
\begin{equation}
    \boldsymbol{C}(\boldsymbol{p}) = \sum_{i \in \mathcal{N}}
    \boldsymbol{c}_i \, \alpha_i \prod_{j=1}^{i-1}(1-\alpha_j),
\end{equation}
where $\alpha_i$ is the product of the 2D-projected Gaussian opacity and
the per-pixel opacity.  Depth rendering is similarly obtained by
alpha-compositing the Gaussian centers along each ray.

\subsection{From NeRF-SLAM to Gaussian-SLAM}

Early radiance-field SLAM systems used NeRF-based implicit
representations~\cite{sucar2021_imap,zhu2022_niceslam,wang2023_coslam,
rosinol2023_nerfslam,sandstrom2023_pointslam}.  While achieving high-quality
rendering, these methods were limited by slow volumetric rendering and
MLP-capacity bottlenecks.  The transition to 3DGS brought three key
advantages: (1)~real-time rasterization ($>$100~FPS), (2)~explicit
primitives enabling direct geometric operations (ICP, nearest-neighbor
queries), and (3)~incremental map expansion without retraining.

\subsection{The Depth-Source Challenge}

In a typical 3DGS-SLAM pipeline, the depth source serves three roles:
\begin{enumerate}
    \item \textbf{Initialization}: Each Gaussian's center $\boldsymbol{\mu}$
    is initialized by back-projecting depth pixels into 3D world coordinates.
    Sparse or inaccurate depth leads to under-reconstruction.
    \item \textbf{Supervision}: Depth rendering loss
    $\mathcal{L}_{\text{depth}} = \|\hat{D} - D\|_1$ constrains the Gaussian
    geometry.  Without reliable depth, Gaussians may converge to incorrect
    positions that overfit training views.
    \item \textbf{Tracking}: In map-centric SLAM (e.g.,
    MonoGS~\cite{matsuki2024monogs}), rendered depth is compared against
    sensor depth to estimate camera pose.  Depth quality directly impacts
    tracking accuracy.
\end{enumerate}

This depth-source dependency is the central axis along which we organize
this survey.  Table~\ref{tab:depth_sources} summarizes the three modalities.

\begin{table}[t]
\centering
\caption{Comparison of depth source modalities for Gaussian Splatting SLAM.}
\label{tab:depth_sources}
\small
\begin{tabular}{lcccc}
\toprule
\textbf{Modality} & \textbf{Metric Scale} & \textbf{Range} & \textbf{Cost} & \textbf{Sparsity} \\
\midrule
LiDAR + RGB + IMU & Yes & Long ($>$100\,m) & High & Sparse \\
iToF + RGB + IMU & Yes & Short ($<$10\,m) & Low & Very Sparse \\
Mono Depth + RGB + IMU & Via IMU & Unlimited & Minimal & Dense (learned) \\
\bottomrule
\end{tabular}
\end{table}

\subsection{Evaluation Metrics}

Gaussian SLAM systems are typically evaluated on:
\begin{itemize}
    \item \textbf{Tracking}: Absolute Trajectory Error (ATE) in cm or m.
    \item \textbf{Rendering}: PSNR (dB), SSIM, LPIPS on held-out views.
    \item \textbf{Geometry}: Depth L1 error, mesh accuracy/completeness
    (cm), Chamfer distance.
    \item \textbf{Efficiency}: FPS, GPU memory (GB), CPU utilization.
\end{itemize}
Common datasets include Replica (indoor), TUM RGB-D, ScanNet, ETH3D, KITTI
(outdoor driving), and the MARS-LVIG dataset for LiDAR-visual-inertial
sequences.


%======================================================================
\section{LiDAR + RGB + IMU Gaussian Splatting SLAM}
\label{sec:lidar}
LiDAR-Inertial-Visual (LIV) fusion has been extensively validated in
classical SLAM~\cite{xu2022fastlio2,zheng2024_fastlivo2,lin2024_r3live},
providing robust metric-scale odometry through tightly-coupled iterated
extended Kalman filters (IESKF) or continuous-time factor-graph optimization.
Integrating 3DGS into this framework creates systems that combine precise
pose estimation with photo-realistic dense mapping.  We categorize the
literature into \emph{tightly-coupled} and \emph{decoupled} approaches
based on whether the Gaussian map feeds back into the odometry.

\subsection{Gaussian-LIC and Gaussian-LIC2}

\textbf{Gaussian-LIC}~\cite{lang2025_gaussianlic} (ICRA 2025) is the first
real-time photo-realistic SLAM system combining 3DGS with
LiDAR-Inertial-Camera fusion.  It uses a continuous-time trajectory
optimization framework for robust pose estimation, initializes Gaussians
from both LiDAR points and triangulated visual SfM points, and models sky
and varying camera exposure for outdoor rendering.  The system is
implemented entirely in C++/CUDA with carefully designed acceleration
strategies.  \textbf{Code}: \url{https://github.com/APRIL-ZJU/Gaussian-LIC}.

\textbf{Gaussian-LIC2}~\cite{lang2026_gaussianlic2} (IJRR 2026) extends this
to a journal version addressing three key issues: (1)~under-reconstruction
in LiDAR-blind areas, solved via a lightweight zero-shot depth completion
network that fuses RGB appearance cues with sparse LiDAR measurements;
(2)~geometric accuracy, enhanced by using sparse but precise LiDAR depths to
supervise Gaussian optimization with CUDA-accelerated strategies;
(3)~bidirectional coupling between odometry and mapping, where photometric
constraints from the incrementally built Gaussian map are fed back into
continuous-time factor graph optimization to improve tracking under LiDAR
degeneration.  A dedicated dataset with ground-truth poses, depths, and
extrapolated trajectories is released for evaluating out-of-sequence novel
view synthesis.  \textbf{Code}: \url{https://xingxingzuo.github.io/gaussian_lic2}.

\subsection{GS-LIVM}

GS-LIVM~\cite{xie2024_gslivm} (ICCV 2025 Workshop) targets real-time
photo-realistic LiDAR-Inertial-Visual mapping for large-scale outdoor
scenes.  Its key innovation is the use of Gaussian Process Regression (GPR)
at the voxel level (Voxel-GPR) to densify sparse and unevenly distributed
LiDAR observations into uniform point grids.  The estimated covariance from
GPR initializes the scale and rotation of 3D Gaussians, accelerating
convergence.  Custom CUDA kernels enable real-time performance.  Evaluated
on outdoor datasets including the Botanic Garden sequence ($\sim$200\,m
trajectory).  \textbf{Code}: \url{https://github.com/xieyusen/GS-LIVM}.

\subsection{GS-LIVO}

GS-LIVO~\cite{hong2025_gslivo} (IEEE T-RO 2025) proposes a global Gaussian
map with hash-indexed voxels organized in a recursive octree, adapting to
different environment scales.  It uses an IESKF for tightly-coupled
multi-sensor fusion and maintains a sliding window of Gaussians to reduce
GPU memory.  Notably, it is the first Gaussian-based SLAM deployable on
embedded systems (NVIDIA Jetson Orin NX), demonstrating real-time performance
on resource-constrained platforms---a critical milestone for robotic
deployment.

\subsection{LVI-GS}

LVI-GS~\cite{zhao2025_lvigs} (IEEE T-IM 2025) introduces a
tightly-coupled LiDAR-Visual-Inertial mapping framework.  Gaussians are
initialized from colorized LiDAR points and optimized via differentiable
rendering.  A pyramid-based training approach learns multi-level features,
and a depth loss derived from LiDAR measurements improves geometric
perception.  Thread management and custom CUDA acceleration enable real-time
operation.  \textbf{Project page}: \url{https://kwanwaipang.github.io/LVI-GS/}.

\subsection{RMGS-SLAM}

RMGS-SLAM~\cite{li2026_rmgs} (arXiv 2026) is a tightly-coupled LIV 3DGS-SLAM
with explicit loop closure.  Its cascaded Gaussian initialization combines
feed-forward predictions with voxel-PCA geometric priors, improving
initialization quality and convergence speed.  Loop closure is performed
directly on the optimized Gaussian map via Gaussian-based GICP registration,
followed by pose-graph optimization.  A new benchmark with
hardware-synchronized solid-state LiDAR, camera, and IMU plus ground truth
is released.  The system achieves a strong balance among real-time
efficiency, localization accuracy, and rendering quality.

\subsection{Real-Time LiDAR Gaussian Splatting SLAM}

Tak et al.~\cite{tak2026_realtimelidar} (arXiv 2026) couple fast G-ICP
registration with spherical rasterization-based dense mapping.  A
covariance-derived geometry score measures local complexity and drives
pruning in planar regions and selective densification in structurally rich
areas, enabling real-time large-scale mapping.

\subsection{LV-GS SLAM}

LV-GS SLAM~\cite{he2026lvgsslam} (Applied Sciences 2026) integrates LiDAR
and visual data for large-scale incremental reconstruction with real-time
tracking.  A LiDAR odometry front-end processes 30\,Hz LiDAR at 15--19\,Hz
tracking.  A novel depth propagation approach initializes 3D Gaussians from
dense depth maps, achieving 9$\times$ faster PSNR convergence versus direct
LiDAR initialization.  Keyframe-based submap management dynamically adjusts
memory allocation to prevent GPU overflow.  Validated on KITTI and a
quadruped robot platform.

\subsection{LIT-GS: LiDAR-Inertial-Thermal}

LIT-GS~\cite{shi2026_litgs} (IROS 2026) extends the LIV paradigm to thermal
imaging, replacing the RGB camera with a long-wave thermal sensor for
illumination-robust mapping.  It injects LiDAR-derived plane geometry as an
explicit constraint in both pose/structure refinement (via weighted
point-to-plane BA) and Gaussian optimization (via a LiDAR-plane-regularized
splatting objective).  This addresses the weak texture and low contrast of
thermal images, which would otherwise corrupt photometric supervision.

\subsection{Summary of LiDAR-Based Systems}

\Cref{tab:lidar_summary} compares the key characteristics of LiDAR+RGB+IMU
Gaussian SLAM systems.  The evolution shows clear trends: (1)~from
offline to real-time operation, (2)~from LiDAR-only initialization to
hybrid LiDAR+depth-completion initialization, (3)~from decoupled
odometry-mapping to bidirectional coupling where the Gaussian map improves
tracking, and (4)~from desktop GPUs to embedded deployment.

\begin{table}[t]
\centering
\caption{LiDAR + RGB + IMU Gaussian Splatting SLAM systems.}
\label{tab:lidar_summary}
\scriptsize
\begin{tabular}{lccccc}
\toprule
\textbf{System} & \textbf{Year} & \textbf{Venue} & \textbf{Real-time} & \textbf{Loop Closure} & \textbf{Code} \\
\midrule
Gaussian-LIC~\cite{lang2025_gaussianlic} & 2025 & ICRA & \checkmark & -- & \checkmark \\
GS-LIVM~\cite{xie2024_gslivm} & 2024 & arXiv & \checkmark & -- & \checkmark \\
GS-LIVO~\cite{hong2025_gslivo} & 2025 & T-RO & \checkmark & -- & -- \\
LVI-GS~\cite{zhao2025_lvigs} & 2025 & T-IM & \checkmark & -- & -- \\
Gaussian-LIC2~\cite{lang2026_gaussianlic2} & 2026 & IJRR & \checkmark & -- & \checkmark \\
RMGS-SLAM~\cite{li2026_rmgs} & 2026 & arXiv & \checkmark & \checkmark & -- \\
RT-LiDAR-GS~\cite{tak2026_realtimelidar} & 2026 & arXiv & \checkmark & -- & -- \\
LV-GS SLAM~\cite{he2026lvgsslam} & 2026 & AppSci & \checkmark & -- & -- \\
LIT-GS~\cite{shi2026_litgs} & 2026 & IROS & \checkmark & -- & -- \\
\bottomrule
\end{tabular}
\end{table}

\textbf{Key insight}: The dominant challenge in LiDAR-based Gaussian SLAM is
\emph{LiDAR sparsity and FoV mismatch with the camera}.  Gaussian-LIC2's
depth completion network and RMGS-SLAM's cascaded initialization are the two
most effective solutions, while GS-LIVM's Voxel-GPR densification provides a
lightweight alternative.  The bidirectional coupling of Gaussian-LIC2
(map$\rightarrow$odometry photometric feedback) is a unique contribution
that demonstrates the Gaussian map is not merely a passive visualization
tool but an active constraint source for tracking.


%======================================================================
\section{iToF Depth + RGB + IMU Gaussian Splatting SLAM}
\label{sec:itof}
Indirect Time-of-Flight (iToF) sensors measure depth by correlating the
phase shift of amplitude-modulated illumination, providing metric depth at
lower cost and power than mechanical LiDAR.  Modern iToF sensors are
ubiquitous in smartphones (e.g., Apple LiDAR with $\leq$576 points) and
emerging consumer devices (as few as 64 points at 200\,mW).  This extreme
sparsity---potentially $8\times8$ depth points per frame---poses a unique
challenge for Gaussian Splatting SLAM, which traditionally requires dense
depth for initialization and supervision.

\subsection{ToF-Splatting}

ToF-Splatting~\cite{conti2025_tofsplatting} (CVPR 2025) is the first 3DGS-based
SLAM pipeline tailored for very sparse ToF input.  Its architecture consists
of three modules:

\begin{enumerate}
    \item \textbf{Tracking Frontend}: A frame-to-model approach that
    minimizes a combined photometric ($\mathcal{L}_{\text{rgb}}$) and
    geometric ($\mathcal{L}_{\text{depth}}$) loss:
    \begin{equation}
        \mathcal{L}_{\text{track}} = \lambda_{\text{track}}
        \mathcal{L}_{\text{rgb}} + (1-\lambda_{\text{track}})
        \mathcal{L}_{\text{depth}}
    \end{equation}
    where $\mathcal{L}_{\text{depth}}$ uses the dense depth predicted by the
    multi-frame integration module (not the raw sparse ToF depth).

    \item \textbf{Multi-Frame Integration (MFI)}: Extends the Depth on Demand
    (DoD) framework~\cite{conti2024_dod} to fuse three complementary cues:
    (a)~extremely sparse ToF depth, (b)~monocular color cues via a Depth
    Anything v2~\cite{yang2024depthanythingv2} encoder, and
    (c)~multi-view geometry from a keyframe buffer.  An outlier rejection
    step compares ToF measurements against a no-ToF depth prediction and
    discards measurements with $\ell_1$ error above a quantile threshold,
    producing filtered sparse depth $\tilde{H}_k$ before final DoD
    prediction.

    \item \textbf{Mapping Backend}: Initializes new Gaussians where rendered
    opacity is low (uncertainty-driven).  The mapping loss combines
    photometric, DSSIM, depth, normal, and isotropy regularization terms:
    \begin{equation}
        \mathcal{L}_{\text{map}} = \lambda_{\text{map}}[\lambda_{\text{visual}}
        \mathcal{L}_{\text{rgb}} + (1-\lambda_{\text{visual}})
        \mathcal{L}_{\text{dssim}}] + (1-\lambda_{\text{map}})
        \mathcal{L}_{\text{depth}} + \lambda_{\text{normals}}
        \mathcal{L}_{\text{normals}} + \lambda_{\text{iso}}
        \mathcal{L}_{\text{iso}}
    \end{equation}
    The isotropy loss $\mathcal{L}_{\text{iso}}$ penalizes elongated
    Gaussians, and a keyframe novelty factor $\nu_k$ based on rendered
    opacity drives keyframe selection.
\end{enumerate}

\textbf{Key results}: On the ZJUL5 dataset (real sparse ToF) and Replica,
ToF-Splatting achieves state-of-the-art tracking and mapping, outperforming
prior sparse-ToF SLAM (ToF-SLAM~\cite{jiang2022quantizedtof}) and achieving
comparable or better results than RGB-D methods even with $\leq$1\% of the
depth points.  On TUM RGB-D, it matches full-density RGB-D Gaussian SLAM.

\textbf{Code}: Project page at
\url{https://andreaconti.github.io/projects/tof_splatting}.

\subsection{iToF-RGB Depth Fusion (Non-SLAM)}

While not SLAM systems per se, several works address the fundamental
iToF+RGB fusion problem that underpins this category:

\textbf{Fang et al.}~\cite{fang2025_iToFRGB} (ICPR 2024) propose a
weakly-supervised framework fusing narrow-FoV iToF depth with wide-FoV RGB
monocular depth.  The learning objective is inter-camera geometric
consistency, supplemented by structure distillation from the passive
sensing modality.  This is the smartphone-targeted counterpart to
ToF-Splatting.

\textbf{Du et al.}~\cite{du2025_iToFRGB} (arXiv 2025) present an iToF-RGB
fusion framework addressing iToF's low spatial resolution, limited FoV, and
structural distortion.  A precise geometric calibration module reprojects
narrow-FoV iToF depth onto the wide-FoV RGB coordinate system, and a
dual-encoder fusion network jointly extracts complementary features guided
by monocular depth priors.

\subsection{ToF-IP: Inertial + ToF for Motion Capture}

ToF-IP~\cite{yao2025_tofip} (NeurIPS 2025) integrates ToF sensors with
sparse IMUs for human motion capture---a different application domain but
demonstrating the complementary nature of ToF (direct distance measurement)
and inertial sensing (indirect acceleration integration).  A Transformer
encoder with Node-Centric Data Integration and Dynamic Spatial Positional
Encoding effectively fuses the two modalities.

\subsection{Summary and Gap Analysis}

The iToF + RGB + IMU Gaussian SLAM category is currently the least explored
of the three, with ToF-Splatting being the sole dedicated 3DGS-SLAM system.
The key gaps are:

\begin{itemize}
    \item \textbf{No IMU integration}: ToF-Splatting uses only ToF + RGB,
    not IMU.  An iToF + RGB + IMU Gaussian SLAM system would combine the
    metric scale from iToF, motion continuity from IMU, and appearance from
    RGB---a natural and unexplored configuration for mobile/AR devices.
    \item \textbf{Outdoor extension}: iToF sensors suffer from ambient light
    interference outdoors.  Extending to outdoor scenes requires robust
    multi-frame integration under varying SNR.
    \item \textbf{Embedded deployment}: iToF sensors target mobile devices,
    but current 3DGS optimization is GPU-intensive.  Lightweight Gaussian
    representations (e.g., LoCoSplat~\cite{wang2026_locosplat}) may enable
    on-device real-time mapping.
\end{itemize}

\begin{table}[t]
\centering
\caption{iToF + RGB + IMU Gaussian Splatting and related systems.}
\label{tab:itof_summary}
\scriptsize
\begin{tabular}{lcccc}
\toprule
\textbf{System} & \textbf{Year} & \textbf{Venue} & \textbf{3DGS SLAM} & \textbf{IMU} \\
\midrule
ToF-Splatting~\cite{conti2025_tofsplatting} & 2025 & CVPR & \checkmark & -- \\
ToF-SLAM~\cite{jiang2022quantizedtof} & 2022 & CVPR-W & -- (NeRF) & -- \\
DoD~\cite{conti2024_dod} & 2024 & ECCV & -- (depth only) & -- \\
iToF-RGB Fusion~\cite{fang2025_iToFRGB} & 2024 & ICPR & -- & -- \\
iToF-RGB Fusion~\cite{du2025_iToFRGB} & 2025 & arXiv & -- & -- \\
ToF-IP~\cite{yao2025_tofip} & 2025 & NeurIPS & -- (MoCap) & \checkmark \\
\bottomrule
\end{tabular}
\end{table}


%======================================================================
\section{Monocular Depth + RGB + IMU Gaussian Splatting SLAM}
\label{sec:monodepth}
When no active depth sensor is available, Gaussian Splatting SLAM must
derive geometric priors from learned models or visual-inertial
structure-from-motion.  This category---monocular depth + RGB + IMU---is the
most hardware-accessible (requiring only a camera and IMU, both ubiquitous
in consumer devices) but the most algorithmically challenging, as monocular
depth is scale-ambiguous, view-dependent, and noisy.

We subdivide this category into three threads: (1)~visual-inertial
Gaussian SLAM with tightly-coupled IMU fusion, (2)~monocular RGB-only
Gaussian SLAM with learned depth priors, and (3)~feed-forward foundation-model-
based Gaussian SLAM.

\subsection{Visual-Inertial Gaussian Splatting SLAM}

\subsubsection{VIGS-SLAM}

VIGS-SLAM~\cite{zhu2026_vigssslam} (arXiv 2025, ETH Zurich) is the
state-of-the-art visual-inertial 3DGS SLAM system.  It tightly couples
visual and inertial energy terms within a unified optimization framework,
jointly refining camera poses, depths, and IMU states (velocity, bias).
Key contributions include:
\begin{itemize}
    \item \textbf{Custom CUDA kernels} for inertial bundle adjustment that
    directly fuse visual and inertial Hessian contributions in a single
    normal-equation solve per iteration---eliminating the linearization gap
    of cascaded approaches like DBA-Fusion.
    \item \textbf{Robust multi-stage IMU initialization} recovering metric
    scale, gravity direction, and IMU parameters.
    \item \textbf{Time-varying bias modeling} (not fixed bias as in prior
    work).
    \item \textbf{Efficient loop closure} with consistent Gaussian map
    updates.
\end{itemize}
Evaluated on five datasets including a self-captured visual-inertial
dataset, VIGS-SLAM consistently outperforms SOTA methods under motion blur,
low texture, and exposure variations.  \textbf{Project page}:
\url{https://visionlab.ethz.ch/vigs-slam}.

\textbf{Elevator-VIGS}~\cite{zhou2026_elevatorvigs} extends VIGS-SLAM with
a transport state that separates elevator motion from robot motion,
enabling tracking through elevator rides---a scenario where the camera and
IMU observe different motions.

\subsubsection{MM3DGS SLAM}

MM3DGS SLAM~\cite{sun2024_mm3dgs} (IROS 2024) is the first system to use
3D Gaussians for map representation with unposed camera images and inertial
measurements.  Its keyframe-based framework incorporates relative pose
transformations from pre-integrated IMU measurements, depth estimates, and
photometric rendering quality.  For monocular input, it treats learned
monocular depth as sensor depth.  Achieves 3$\times$ tracking improvement
over prior 3DGS SLAM SOTA.

\subsubsection{CGS-SLAM}

CGS-SLAM~\cite{deambrogi2026_cgssslam} (arXiv 2026) enables multi-agent
3DGS SLAM using only RGB and inertial data.  Each agent performs local
tracking with inertial data as a motion prior and reconstructs a scaled map
using a \emph{metric} monocular depth estimator (Depth Pro).  Keyframe
encodings are shared among agents for dynamic keyframing in overlap
regions, and a central server aligns submaps using VGGT~\cite{wang2025_vggt}
as a view alignment model.  This is the collaborative counterpart to
VIGS-SLAM, targeting GNSS-denied tactical scenarios with consumer
smartphones.

\subsubsection{VIM-GS}

VIM-GS~\cite{zhang2026_vimgs} (arXiv 2025) is a visual-inertial monocular
Gaussian Splatting framework for large scenes.  It addresses the depth
challenge by fusing accurate but sparse depth from visual-inertial SfM with
dense but coarse depth from large foundation models (LFMs).  An
object-segmented depth propagation algorithm bridges sparse and dense
depth, and a dynamic depth refinement module handles crippled SfM depth of
dynamic objects.  This is the large-scale outdoor counterpart to indoor
VIGS-SLAM.

\subsection{Monocular RGB-Only Gaussian SLAM with Learned Depth}

\subsubsection{MonoGS (Gaussian Splatting SLAM)}

MonoGS~\cite{matsuki2024monogs} (CVPR 2024) is the first application of 3D
Gaussian Splatting in monocular SLAM.  It uses 3DGS as the sole scene
representation and introduces geometric regularization to address
ambiguities in incremental reconstruction from monocular input.  The system
demonstrates near real-time operation and achieves SOTA results in novel
view synthesis and trajectory estimation, including reconstruction of tiny
and transparent objects.  \textbf{Code}:
\url{https://github.com/muskie82/MonoGS}.

\subsubsection{HI-SLAM2}

HI-SLAM2~\cite{zhang2026_hislam2} (IEEE T-RO 2025) achieves fast and
accurate monocular scene reconstruction using only RGB input.  Its key idea
is combining easy-to-obtain monocular geometry priors (depth + normals from
OmniData) with learning-based dense SLAM, using 3DGS as the map
representation.  A grid-based scale alignment strategy addresses scale
distortions in monocular depth priors.  Upon loop closure, efficient
pose-graph BA with instant Gaussian deformation ensures global consistency.
Surpasses even RGB-D methods on Replica, ScanNet, and ScanNet++.  \textbf{Code}:
\url{https://hi-slam2.github.io/}.

\subsubsection{Splat-SLAM}

Splat-SLAM~\cite{sandstrom2025_splatslam} (CVPR 2025) is the first RGB-only
SLAM system with globally optimized tracking, dynamically adapting the
Gaussian map to keyframe pose and depth updates.  It incorporates
Disparity, Scale and Pose Optimization (DSPO) for bundle adjustment,
jointly optimizing pose, depth, and monocular depth scale.

\subsubsection{GaussianFlow SLAM}

GaussianFlow SLAM~\cite{seo2026_gaussianflow} (IEEE RA-L 2026) leverages
optical flow as a geometry-aware cue to guide the optimization of both the
3DGS and camera poses.  By aligning the projected motion of Gaussians
(``GaussianFlow'') with dense optical flow, it continuously corrects 3DGS
geometry even without direct depth measurements.  Normalized error-based
densification and pruning modules refine inactive Gaussians.  \textbf{Code}:
\url{https://github.com/url-kaist/gaussianflow-slam}.

\subsubsection{MyGO-Splat}

MyGO-Splat~\cite{zhu2026_mygosplat} (arXiv 2026) is a closed-loop Gaussian
SLAM framework that analytically rasterizes Gaussian primitives into
pixel-wise depth and surface normals, allowing the map to actively supervise
camera pose optimization.  Achieves performance comparable to RGB-D methods
while using only monocular input.

\subsubsection{WildGS-SLAM}

WildGS-SLAM~\cite{zheng2025_wildgssslam} (CVPR 2025) addresses monocular
Gaussian SLAM in dynamic environments.  An uncertainty map, predicted by a
shallow MLP given DINOv2 features, guides dynamic object removal during
tracking (weighted DBA) and mapping (weighted rendering loss).  Monocular
metric depth facilitates pose estimation.  A new Wild-SLAM dataset with
diverse indoor/outdoor dynamic scenes is contributed.

\subsubsection{GeoGS-SLAM}

GeoGS-SLAM~\cite{zhou2026_geogs} (arXiv 2026) introduces geometry-only
Gaussian Splatting (GeoGS) that directly reconstructs scene geometry, and
builds a dense visual SLAM system upon it.  It outperforms SOTA methods in
online mapping efficiency and geometric reconstruction quality.

\subsubsection{MonoEM-GS}

MonoEM-GS~\cite{kruzhkov2026_monoemgs} (arXiv 2026) integrates feed-forward
geometric foundation model predictions into a global Gaussian Splatting
representation via an Expectation-Maximization formulation that explicitly
addresses view-dependent geometry inconsistencies.  ICP-based alignment
provides monocular pose estimation.  Parameterizes Gaussians with
multi-modal features for in-place open-set segmentation.

\subsection{Feed-Forward Foundation-Model-Based Gaussian SLAM}

\subsubsection{VGGT-GS SLAM}

VGGT-GS SLAM~\cite{han2026_vggtgssslam} (arXiv 2026) is a monocular 3DGS
SLAM system for uncalibrated videos.  Starting from feed-forward VGGT
pose and depth priors, it performs submap differentiable bundle adjustment
that jointly refines camera poses, a 3D Gaussian map, and
submap-shared intrinsics and radial-tangential distortion through analytic
calibration Jacobians.  Gaussian-native alignment (GNA) provides
camera-anchored scale refinement between submaps and loop-closure
verification.

\subsection{Summary of Monocular-Depth-Based Systems}

\Cref{tab:mono_summary} compares the key monocular-depth Gaussian SLAM
systems.  The field has progressed rapidly from the initial MonoGS
(RGB-only, near real-time) to VIGS-SLAM (tightly-coupled VI, real-time, loop
closure) and foundation-model-based systems (VGGT-GS SLAM, uncalibrated).

\begin{table}[t]
\centering
\caption{Monocular Depth + RGB + IMU Gaussian Splatting SLAM systems.}
\label{tab:mono_summary}
\scriptsize
\begin{tabular}{lcccccc}
\toprule
\textbf{System} & \textbf{Year} & \textbf{Venue} & \textbf{IMU} & \textbf{Depth Prior} & \textbf{Loop} & \textbf{Code} \\
\midrule
MonoGS~\cite{matsuki2024monogs} & 2024 & CVPR & -- & None & \checkmark & \checkmark \\
MM3DGS~\cite{sun2024_mm3dgs} & 2024 & IROS & \checkmark & Mono depth & -- & -- \\
HI-SLAM2~\cite{zhang2026_hislam2} & 2025 & T-RO & -- & OmniData & \checkmark & \checkmark \\
Splat-SLAM~\cite{sandstrom2025_splatslam} & 2025 & CVPR & -- & Mono depth & \checkmark & -- \\
WildGS-SLAM~\cite{zheng2025_wildgssslam} & 2025 & CVPR & -- & Metric depth & -- & -- \\
VIM-GS~\cite{zhang2026_vimgs} & 2025 & arXiv & \checkmark & LFM + SfM & -- & -- \\
VIGS-SLAM~\cite{zhu2026_vigssslam} & 2025 & arXiv & \checkmark & DROID-SLAM & \checkmark & -- \\
GaussianFlow~\cite{seo2026_gaussianflow} & 2026 & RA-L & -- & Optical flow & -- & \checkmark \\
GeoGS-SLAM~\cite{zhou2026_geogs} & 2026 & arXiv & -- & Learned & -- & -- \\
MyGO-Splat~\cite{zhu2026_mygosplat} & 2026 & arXiv & -- & Analytic & -- & -- \\
MonoEM-GS~\cite{kruzhkov2026_monoemgs} & 2026 & arXiv & -- & Foundation & -- & -- \\
VGGT-GS SLAM~\cite{han2026_vggtgssslam} & 2026 & arXiv & -- & VGGT & \checkmark & -- \\
CGS-SLAM~\cite{deambrogi2026_cgssslam} & 2026 & arXiv & \checkmark & Depth Pro & -- & -- \\
\bottomrule
\end{tabular}
\end{table}

\textbf{Key insight}: The fundamental challenge in monocular-depth Gaussian
SLAM is \emph{scale ambiguity}.  Three approaches address this:
(a)~IMU pre-integration for metric scale (VIGS-SLAM, MM3DGS, CGS-SLAM),
(b)~grid-based scale alignment of monocular depth priors (HI-SLAM2), and
(c)~metric monocular depth estimation models (Depth Pro in CGS-SLAM,
Metric3D in MAGS-SLAM).  The IMU-based approach is most robust but requires
hardware; the metric depth model approach is hardware-free but relies on
generalization of the depth network.


%======================================================================
\section{Cross-Category Comparison and Discussion}
\label{sec:comparison}
\subsection{Accuracy and Robustness Trade-offs}

The three depth-source modalities occupy distinct positions in the
accuracy--cost--robustness trade-off space:

\begin{itemize}
    \item \textbf{LiDAR + RGB + IMU}: Highest tracking accuracy and
    geometric fidelity, especially outdoors and under aggressive motion.
    LiDAR provides direct metric geometry that anchors the Gaussian map.
    However, LiDAR hardware is expensive ($$\,$$1k--$$\,$$10k+), bulky, and
    power-hungry.  Systems like Gaussian-LIC2 and RMGS-SLAM achieve
    sub-decimeter ATE in outdoor scenes where visual-only methods fail.

    \item \textbf{iToF + RGB + IMU}: Moderate accuracy, limited by iToF's
    short range ($<$10\,m) and ambient light sensitivity.  The extreme
    sparsity ($\leq$64 points) requires sophisticated depth completion.
    ToF-Splatting demonstrates that multi-frame integration can bridge the
    gap to RGB-D quality, but the approach remains indoor-focused.
    Hardware cost is minimal ($$\,$$5--$$\,$$50), targeting smartphones
    and AR/VR headsets.

    \item \textbf{Monocular Depth + RGB + IMU}: Accuracy depends heavily on
    the depth prior quality and IMU fusion strategy.  VIGS-SLAM achieves
    competitive results by tightly coupling IMU with learned dense
    correspondences.  HI-SLAM2 surpasses RGB-D methods on indoor benchmarks
    through monocular geometry priors + scale alignment.  However, outdoor
    large-scale scenes remain challenging due to depth ambiguity at range.
    Hardware cost is essentially zero (camera + IMU already in every
    smartphone).
\end{itemize}

\subsection{Depth Completion and Densification Strategies}

A common thread across all three categories is the need to \emph{densify}
the raw depth source before initializing Gaussians.  \Cref{tab:depth_comp}
summarizes the strategies:

\begin{table}[t]
\centering
\caption{Depth completion/densification strategies across modalities.}
\label{tab:depth_comp}
\scriptsize
\begin{tabular}{lll}
\toprule
\textbf{Modality} & \textbf{Strategy} & \textbf{Representative} \\
\midrule
LiDAR & Zero-shot depth completion network & Gaussian-LIC2 \\
LiDAR & Voxel-GPR interpolation & GS-LIVM \\
LiDAR & Voxel-PCA geometric priors & RMGS-SLAM \\
LiDAR & Depth propagation from dense maps & LV-GS SLAM \\
\midrule
iToF & Multi-frame integration (DoD) & ToF-Splatting \\
iToF & Monocular depth prior fusion & Fang et al. \\
\midrule
Mono & Foundation model depth (Depth Pro) & CGS-SLAM \\
Mono & OmniData depth + scale grid & HI-SLAM2 \\
Mono & DROID-SLAM dense disparity & VIGS-SLAM \\
Mono & Optical flow as geometry cue & GaussianFlow \\
Mono & VGGT feed-forward pointmaps & VGGT-GS SLAM \\
Mono & SfM + LFM fusion & VIM-GS \\
\bottomrule
\end{tabular}
\end{table}

\subsection{Real-Time Performance}

Real-time operation is a critical differentiator.  The key bottleneck is
Gaussian optimization, which involves differentiable rasterization and
gradient backpropagation at each keyframe.  Strategies to achieve real-time
performance include:

\begin{itemize}
    \item \textbf{Sliding-window optimization}: GS-LIVO optimizes only
    Gaussians within a temporal sliding window, reducing GPU memory and
    computation.
    \item \textbf{Custom CUDA kernels}: Gaussian-LIC, Gaussian-LIC2, LVI-GS,
    and VIGS-SLAM implement custom CUDA kernels for rasterization and
    optimization, avoiding Python overhead.
    \item \textbf{Submap management}: RMGS-SLAM and LV-GS SLAM use
    keyframe-based submaps with dynamic memory allocation.
    \item \textbf{Parallel tracking and mapping}: Most systems separate
    odometry and Gaussian mapping into parallel threads.
    \item \textbf{Embedded deployment}: GS-LIVO (Jetson Orin NX) and
    GS-LIVM (custom CUDA) demonstrate feasibility on resource-constrained
    platforms.
\end{itemize}

\subsection{Open-Source Code Availability}

Open-source code is critical for reproducibility and community adoption.
\Cref{tab:code} lists the systems with confirmed code releases:

\begin{table}[t]
\centering
\caption{Open-source code availability.}
\label{tab:code}
\scriptsize
\begin{tabular}{lcc}
\toprule
\textbf{System} & \textbf{GitHub URL} & \textbf{Category} \\
\midrule
Gaussian-LIC & APRIL-ZJU/Gaussian-LIC & LiDAR \\
GS-LIVM & xieyusen/GS-LIVM & LiDAR \\
Gaussian-LIC2 & xingxingzuo/gaussian\_lic2 & LiDAR \\
MonoGS & muskie82/MonoGS & Monocular \\
HI-SLAM2 & hi-slam2.github.io & Monocular \\
GaussianFlow & url-kaist/gaussianflow-slam & Monocular \\
ToF-Splatting & andreaconti.github.io & iToF \\
\bottomrule
\end{tabular}
\end{table}

\subsection{Loop Closure and Global Consistency}

Loop closure is essential for large-scale SLAM but is often omitted in
Gaussian SLAM systems to preserve runtime.  Only a few systems implement
explicit loop closure:

\begin{itemize}
    \item \textbf{RMGS-SLAM}: Gaussian-based GICP registration on the
    optimized map, followed by pose-graph optimization.
    \item \textbf{VIGS-SLAM}: Optical-flow-based loop detection with
    consistent Gaussian map updates after pose correction.
    \item \textbf{HI-SLAM2}: Sim(3) pose-graph BA with instant Gaussian
    deformation based on anchored keyframe updates.
    \item \textbf{VGGT-GS SLAM}: Gaussian-native alignment for
    loop-candidate verification before global Sim(3) pose-graph
    optimization.
\end{itemize}

The challenge is that after pose correction, the Gaussian map must be
updated to reflect the new keyframe poses---either by re-optimization
(expensive) or by explicit Gaussian deformation (efficient, as in HI-SLAM2).


%======================================================================
\section{Open Challenges and Future Directions}
\label{sec:future}
\subsection{Unified Depth Fusion Frameworks}

Current systems are specialized for a single depth modality.  A unified
framework that can adaptively fuse whatever depth sources are
available---LiDAR, iToF, monocular depth, or combinations thereof---would
maximize flexibility across hardware platforms.  The depth completion
network of Gaussian-LIC2, which fuses RGB appearance with sparse LiDAR,
hints at this direction but remains LiDAR-specific.

\subsection{Learned Front-Ends for Gaussian SLAM}

Most Gaussian SLAM systems use classical odometry front-ends (IESKF,
DROID-SLAM, ORB-SLAM3).  Recent work on streaming multi-camera odometry
with frozen 3D foundation models (StreamRig~\cite{wei2026_streamrig}) and
learned visual-inertial dense BA (VIGS-SLAM's custom CUDA kernels)
suggests that learned front-ends can provide more robust correspondences
for Gaussian initialization, especially under challenging conditions.

\subsection{Memory-Bounded Gaussian Mapping}

Large-scale outdoor mapping with 3DGS generates millions of Gaussians,
exceeding GPU memory.  Current solutions include sliding-window
optimization (GS-LIVO), submap management (RMGS-SLAM, LV-GS SLAM), and
budgeted LOD (Budgeted-GS~\cite{wang2026_budgetedgs}).  A principled
memory-bounded Gaussian allocation strategy---driven by uncertainty
(BayesianGS-SLAM~\cite{kang2026_bayesiangssslam}) or semantic
importance---remains an open problem.

\subsection{Dynamic Scene Handling}

Most Gaussian SLAM systems assume static scenes.  WildGS-SLAM addresses
dynamic distractors via uncertainty estimation, and DynActiveGS decomposes
uncertainty into structural vs.\ motion-induced fields.  However, a
complete dynamic-scene Gaussian SLAM with online moving-object tracking,
map segmentation, and consistent static-scene reconstruction is still
lacking.

\subsection{Edge and Mobile Deployment}

The ultimate goal of multimodal Gaussian SLAM is deployment on
resource-constrained platforms: drones, AR headsets, smartphones.  GS-LIVO's
Jetson deployment is a milestone, but real-time 3DGS on mobile-class GPUs
requires fundamentally lighter representations.  The local-primitive
observation of LoCoSplat~\cite{wang2026_locosplat} (0.14M-param MLP replacing
heavy 3D networks) and the budgeted LOD of Budgeted-GS point toward
practical on-device Gaussian SLAM.

\subsection{Evaluation Standardization}

The field lacks standardized benchmarks for multimodal Gaussian SLAM.
LiDAR-based systems evaluate on MARS-LVIG and KITTI; monocular systems on
Replica, TUM, and ScanNet; iToF systems on ZJUL5.  A unified benchmark with
synchronized multi-sensor data (LiDAR, iToF, RGB, IMU, ground truth) across
indoor/outdoor/dynamic scenes would enable fair cross-category comparison.
Gaussian-LIC2's dataset and RMGS-SLAM's looped outdoor benchmark are steps
in this direction.

\subsection{Bidirectional Map--Odometry Coupling}

Gaussian-LIC2 pioneered feeding photometric constraints from the Gaussian
map back into the odometry, demonstrating improved tracking under LiDAR
degeneration.  This bidirectional coupling---where the map is not merely a
passive visualization but an active constraint source---is a promising but
underexplored direction.  Extending it to visual-inertial systems (where the
Gaussian map could provide depth supervision for the tracking front-end) and
to iToF systems (where the Gaussian map could guide ToF depth completion)
could yield significant robustness gains.


%======================================================================
\section{Conclusion}
\label{sec:conclusion}
This survey has reviewed the state of the art in multimodal Gaussian
Splatting SLAM along three depth-source axes: LiDAR+RGB+IMU,
iToF+RGB+IMU, and monocular depth+RGB+IMU.  The key findings are:

\begin{enumerate}
    \item \textbf{LiDAR-based systems} (Gaussian-LIC2, RMGS-SLAM, GS-LIVO,
    etc.) achieve the highest accuracy and outdoor robustness through
    direct metric geometry, but at high hardware cost.  The field is
    converging on depth completion networks to address LiDAR sparsity and
    bidirectional map--odometry coupling for improved tracking.

    \item \textbf{iToF-based systems} (ToF-Splatting) are in their infancy
    but demonstrate that even $8\times8$ sparse ToF depth, fused with
    multi-view geometry and monocular cues, can achieve competitive SLAM
    performance---a critical enabler for mobile and AR/VR deployment.
    The integration of IMU into iToF Gaussian SLAM remains an open
    opportunity.

    \item \textbf{Monocular-depth-based systems} (VIGS-SLAM, HI-SLAM2,
    VGGT-GS SLAM, etc.) are the most accessible and rapidly evolving,
    leveraging foundation model depth priors and tight visual-inertial
    coupling to approach RGB-D quality without active depth sensors.
    Scale ambiguity remains the core challenge, addressed by IMU fusion,
    metric depth models, or grid-based scale alignment.
\end{enumerate}

The unifying trend across all three categories is the shift from
\emph{passive mapping} (Gaussians as visualization) to \emph{active
mapping} (Gaussians as geometric constraints for tracking, depth
completion, and loop closure).  As 3DGS representations become lighter
and foundation model priors become more accurate, we expect multimodal
Gaussian Splatting SLAM to transition from research prototype to
practical deployment on consumer devices.


\bibliographystyle{IEEEtran}
\bibliography{survey_refs}

\appendix
\section{Paper Inventory}
\label{sec:inventory}
\Cref{tab:full_inventory} lists all papers surveyed in this work, organized
by category.  PDFs for papers with arXiv preprints are saved in
\texttt{latex/papers/} following the naming convention
\texttt{FirstauthorLastName<Year>\_<ShortName>.pdf}.

\begin{table}[h]
\centering
\caption{Complete paper inventory by category.}
\label{tab:full_inventory}
\scriptsize
\begin{tabular}{p{4cm}p{2cm}p{2cm}p{5cm}}
\toprule
\textbf{Paper} & \textbf{Year} & \textbf{Venue} & \textbf{Key Contribution} \\
\midrule
\multicolumn{4}{l}{\textbf{Category 1: LiDAR + RGB + IMU}} \\
\midrule
Gaussian-LIC~\cite{lang2025_gaussianlic} & 2025 & ICRA & First real-time LIV GS-SLAM \\
Gaussian-LIC2~\cite{lang2026_gaussianlic2} & 2026 & IJRR & Depth completion + bidirectional coupling \\
GS-LIVM~\cite{xie2024_gslivm} & 2024 & ICCV-W & Voxel-GPR densification \\
GS-LIVO~\cite{hong2025_gslivo} & 2025 & T-RO & Embedded (Jetson) deployment \\
LVI-GS~\cite{zhao2025_lvigs} & 2025 & T-IM & Pyramid training + depth loss \\
RMGS-SLAM~\cite{li2026_rmgs} & 2026 & arXiv & Cascaded init + loop closure \\
RT-LiDAR-GS~\cite{tak2026_realtimelidar} & 2026 & arXiv & G-ICP + spherical rasterization \\
LV-GS SLAM~\cite{he2026lvgsslam} & 2026 & AppSci & Depth propagation + submap mgmt \\
LIT-GS~\cite{shi2026_litgs} & 2026 & IROS & LiDAR-inertial-thermal \\
\midrule
\multicolumn{4}{l}{\textbf{Category 2: iToF + RGB + IMU}} \\
\midrule
ToF-Splatting~\cite{conti2025_tofsplatting} & 2025 & CVPR & First sparse-ToF 3DGS-SLAM \\
DoD~\cite{conti2024_dod} & 2024 & ECCV & Multi-frame depth integration \\
iToF-RGB Fusion~\cite{fang2025_iToFRGB} & 2024 & ICPR & Weakly-supervised iToF+RGB \\
iToF-RGB Fusion~\cite{du2025_iToFRGB} & 2025 & arXiv & Monocular-aided iToF+RGB \\
ToF-IP~\cite{yao2025_tofip} & 2025 & NeurIPS & ToF+IMU for MoCap \\
\midrule
\multicolumn{4}{l}{\textbf{Category 3: Monocular Depth + RGB + IMU}} \\
\midrule
MonoGS~\cite{matsuki2024monogs} & 2024 & CVPR & First monocular GS-SLAM \\
MM3DGS~\cite{sun2024_mm3dgs} & 2024 & IROS & Multi-modal VI GS-SLAM \\
HI-SLAM2~\cite{zhang2026_hislam2} & 2025 & T-RO & Geometry-aware mono GS-SLAM \\
Splat-SLAM~\cite{sandstrom2025_splatslam} & 2025 & CVPR & Globally optimized RGB-only \\
WildGS-SLAM~\cite{zheng2025_wildgssslam} & 2025 & CVPR & Dynamic-scene mono GS-SLAM \\
VIM-GS~\cite{zhang2026_vimgs} & 2025 & arXiv & VI mono GS for large scenes \\
VIGS-SLAM~\cite{zhu2026_vigssslam} & 2025 & arXiv & Tightly-coupled VI GS-SLAM \\
GaussianFlow~\cite{seo2026_gaussianflow} & 2026 & RA-L & Optical-flow-guided mono GS-SLAM \\
GeoGS-SLAM~\cite{zhou2026_geogs} & 2026 & arXiv & Geometry-only GS dense SLAM \\
MyGO-Splat~\cite{zhu2026_mygosplat} & 2026 & arXiv & Closed-loop geometric feedback \\
MonoEM-GS~\cite{kruzhkov2026_monoemgs} & 2026 & arXiv & EM-based mono GS-SLAM \\
VGGT-GS SLAM~\cite{han2026_vggtgssslam} & 2026 & arXiv & Uncalibrated VGGT-prior GS-SLAM \\
CGS-SLAM~\cite{deambrogi2026_cgssslam} & 2026 & arXiv & Collaborative VI GS-SLAM \\
Elevator-VIGS~\cite{zhou2026_elevatorvigs} & 2026 & arXiv & Elevator-robust VI GS-SLAM \\
\midrule
\multicolumn{4}{l}{\textbf{Foundations}} \\
\midrule
3DGS~\cite{kerbl2023_3dgs} & 2023 & ACM TOG & 3D Gaussian Splatting \\
FAST-LIO2~\cite{xu2022fastlio2} & 2022 & T-RO & LiDAR-inertial odometry \\
FAST-LIVO2~\cite{zheng2024_fastlivo2} & 2024 & T-RO & LIVO odometry \\
R3LIVE~\cite{lin2024_r3live} & 2024 & ICRA & RGB-colored LIVO \\
VGGT~\cite{wang2025_vggt} & 2025 & CVPR & Feed-forward 3D vision FM \\
Depth Anything v2~\cite{yang2024depthanythingv2} & 2024 & arXiv & Monocular depth FM \\
ORB-SLAM3~\cite{campos2021orbslam3} & 2021 & T-RO & VI SLAM \\
VINS-Mono~\cite{qin2018vinsmono} & 2018 & T-RO & Visual-inertial SLAM \\
\bottomrule
\end{tabular}
\end{table}


\end{document}
